\documentclass[UTF8,11pt,a4paper]{ctexart}
\newcommand{\DocTitle}{小组会议议程}
\newcommand{\DocSubtitle}{2026 年 11 月 23 日 · 线上会议}
% Shared layout for the CW2 companion documents.
\usepackage[margin=2.1cm,headheight=16pt]{geometry}
\usepackage{amsmath,amssymb,booktabs,array,tabularx,ragged2e}
\usepackage{enumitem,xcolor,fancyhdr,listings,tcolorbox}
\tcbuselibrary{breakable}
\usepackage{hyperref}
\usepackage{bookmark}

\newcommand{\DocVersion}{3.0}
\newcommand{\EngineVersion}{1.1}
\newcommand{\DocDate}{2026-11-23}
\definecolor{ink}{HTML}{1F2937}
\definecolor{accent}{HTML}{92400E}
\definecolor{muted}{HTML}{4B5563}
\definecolor{linkink}{HTML}{174678}
\definecolor{codebg}{HTML}{F5F7FA}
\hypersetup{unicode,colorlinks=true,linkcolor=linkink,urlcolor=linkink,
  filecolor=linkink,pdftitle={SOE104 CW2 \DocTitle\space v\DocVersion},
  pdfauthor={Group A--E}}
\ctexset{section={format={\Large\bfseries\color{accent}}},
  subsection={format={\large\bfseries\color{ink}}}}
\pagestyle{fancy}
\fancyhf{}
\fancyhead[L]{\small\color{muted}SOE104 CW2 · Oregon Trail}
\fancyhead[R]{\small\color{muted}\DocTitle\ v\DocVersion}
\fancyfoot[C]{\small\color{muted}\thepage}
\renewcommand{\headrulewidth}{0.4pt}
\setlength{\parskip}{0.35em}
\setlength{\emergencystretch}{1.5em}
\renewcommand{\arraystretch}{1.23}
\setlist{itemsep=3pt,topsep=3pt,parsep=0pt,leftmargin=1.7em}
\newcolumntype{Y}{>{\RaggedRight\arraybackslash}X}
\newcolumntype{L}[1]{>{\RaggedRight\arraybackslash}p{#1}}
\newtcolorbox{keybox}[1]{breakable,title={#1},colback=orange!4,
  colframe=accent,boxrule=0.55pt,arc=2pt,fonttitle=\bfseries}
\newtcolorbox{warnbox}[1]{breakable,title={#1},colback=red!3,
  colframe=red!45!black,boxrule=0.55pt,arc=2pt,fonttitle=\bfseries}
\lstset{basicstyle=\small\ttfamily,backgroundcolor=\color{codebg},
  xleftmargin=0.9em,frame=l,rulecolor=\color{linkink!45},
  breaklines=true,columns=fullflexible,keepspaces=true,showstringspaces=false,
  aboveskip=0.65em,belowskip=0.65em}
\newcommand{\code}[1]{\texttt{#1}}
\newcommand{\source}[1]{\textup{\small\textcolor{muted}{[#1]}}}
\newcommand{\docref}[2]{\href{./#1.pdf}{#2}}
\renewcommand{\maketitle}{%
  \begin{center}
    {\LARGE\bfseries CW2 \DocTitle\par}
    \vspace{0.35em}{\large\DocSubtitle\par}
    \vspace{0.75em}{\small\color{muted}Group A--E · 整理：A
      \quad\DocDate\quad 文档 v\DocVersion\par}
    \vspace{0.2em}{\small\color{muted}Engine / Marking v\EngineVersion\par}
  \end{center}
  \vspace{0.4em}}


\begin{document}
\maketitle
\thispagestyle{fancy}

\begin{keybox}{11 月 23 日会议结束前要定下的事}
确认 A--E 的分工、首个快照契约、仓库访问与第一次集成时间。
会前看\docref{05_proposal}{《05 分工与进度提案》}的职责表，以及
\docref{03_plan}{《03 项目实施方案》}的模块与数据流。
\end{keybox}

\section{线上会议，30 分钟}

\begin{center}\small
\begin{tabularx}{\linewidth}{@{}L{2.05cm}L{3.0cm}Y@{}}
\toprule
用时 & 环节 & 当场确认 \\
\midrule
0--2 分钟 & 开场 & A 说明本次目标，课程截止 12-13 23:59（中国时间）。 \\
2--7 分钟 & 目标和规则 & 全员各举一个入口或规则边界。用 01、02 和 04 核对 1.1 口径。 \\
7--12 分钟 & 基线与接口 & A、B、D 确认 \code{Plans} 哪些字段是快照，哪些是投影，缺快照怎么处理。 \\
12--20 分钟 & 分工 & A--E 说明自己交付什么、依赖谁和如何验收。记录需结对的部分。 \\
20--25 分钟 & 版本与合规 & 确认仓库访问和合并方式。约定平台版本记录、回放分工与 AI 申报。 \\
25--30 分钟 & 回顾与行动 & A 复述责任人和时间。未定项标明负责人，放入会后记录。 \\
\bottomrule
\end{tabularx}
\end{center}

\section{会后记录}

\begin{center}\small
\begin{tabularx}{\linewidth}{@{}L{3.55cm}L{2.15cm}Y@{}}
\toprule
记录项 & 谁确认 & 完成标志 \\
\midrule
访问与分工 & A、全员 & 组员可访问现有仓库，确认负责范围与替代方案。 \\
快照接口 & A、B、D & 记录回合号、事件时预测、回退和两三个边界用例。 \\
集成与平台基线 & A、E & 按 05 的目标日期推进，记录版本和三个种子的回放。 \\
\bottomrule
\end{tabularx}
\end{center}

\noindent 口粮阈值、阵容和姿态取舍留给配对回放。
操作步骤见\docref{06_git_handbook}{《06 Git 协作手册》}与
\docref{07_ide_setup}{《07 IDEA 配置》}。

\clearpage
\section{主持人提示与记录样例}

主持人可以用下面的问题确认理解。答复不一致时，先把分歧和要核对的资料记下来，
会后由对应负责人补上结论。

\begin{center}\small
\begin{tabularx}{\linewidth}{@{}L{2.2cm}Y Y@{}}
\toprule
对应时间 & 可以问什么 & 听到什么就可以往下走 \\
\midrule
2--7 分钟 &
策略、事件处理和成员动作分别在什么时候发生？事件扣粮会不会让本回合重新挨饿？ &
能说出先定配速和口粮、再行进喂食、再处理事件和成员行动。
扣粮发生在喂食之后，会影响后续补给。 \\
7--12 分钟 &
给 D 的 \code{Plans} 只有回合初的粮食够不够？事件回合号对不上时怎么办？ &
B、D 确认实际读取值、事件时预测值和缺记录时的回退。
能说出一例，例如喂食会让事件时的粮食变少。 \\
12--20 分钟 &
请每个人说一下：我交什么、要等谁提供什么、别人怎么确认我做完了？ &
A--E 各有一项可核对的交付物。需要结对或互相约定接口的地方有负责人。 \\
20--25 分钟 &
首个需要提交平台的版本怎么记下？谁看回放，谁收集提交流程需要的记录？ &
版本编号、三个种子的记录方式和合规自查有人负责。 \\
\bottomrule
\end{tabularx}
\end{center}

\begin{keybox}{把会后任务写成一条可核对的记录}
例如「任务：确定 \code{Plans} 回合号和快照的读取条件。负责人：A、B、D。
完成时间：会上约定。完成标志：一份接口示例和一个缺快照时的预期结果。」
实际日期和责任人按会上决定填写，更新到\docref{05_proposal}{《05 分工与进度提案》}。
\end{keybox}

\noindent 阈值和采购比例若当场没有证据，记录为回放试验。
例如记录「用同样三个种子比较四头牛与五头牛」，并指定谁汇总到达、失牛和粮食数据。

\end{document}
